% GENERATED FILE. Edit canonical lessons, then run npm run generate:books.
% canonical-source-hash: fnv1a64:4856b4a1aa84f586
% canonical-lessons: IT-C410-guasto, IT-C410-guastarsi, IT-C410-carro-attrezzi, IT-C410-rimorchiare, IT-C410-officina

\chapter{Il guasto}
\label{ch:it-il-guasto}
\hlchaptermodality{410}

\begin{chapteropening}
\textbf{By the end of this chapter:} \emph{I can deal with a breakdown: the fault, breaking down, the tow truck, towing and the garage.}

Complete the last lesson of chapter 410: I can deal with a breakdown: the fault, breaking down, the tow truck, towing and the garage.
\end{chapteropening}

\section[il guasto]{il guasto — the breakdown, the fault}

\label{lesson:IT-C410-guasto}



\begin{quote}
Before the new one: say the Italian for the jack, then the Italian for to pump up.
\end{quote}

\subsection*{You'll want to know: il guasto}
\textbf{il guasto} — \textquotedblleft{}the breakdown, the fault\textquotedblright{}.

\begin{quote}
C'è un guasto al motore.
\end{quote}

\begin{quote}
\emph{There is a fault in the engine.}
\end{quote}

As an adjective, \emph{guasto} means out of order: \emph{il telefono è guasto}.

\begin{cousinweb}
\emph{guasto} is built on \emph{guastare} 'to spoil', which is inherited from Latin \emph{vāstāre} 'to lay waste' (EtymDB).
\end{cousinweb}

\begin{grammarlens}[title={What you've built}]
One more word for telling, explaining and describing at B1.
\end{grammarlens}

\subsection*{Your turn}
\begin{itemize}
\raggedright
  \item \emph{Say it:} \emph{il guasto}
  \item \emph{Say it:} \emph{il guasto}, once more
  \item \emph{Say it:} a sentence of your own with \emph{il guasto}
  \item \emph{Recall:} say the Italian for the jack, then the Italian for to pump up, then say \emph{il guasto} again
\end{itemize}

\begin{tcolorbox}[breakable,colback=teal!4,colframe=teal!35!black,title={Before you move on}]
What does il guasto mean? (\textquotedblleft{}The breakdown, the fault\textquotedblright{}.) Say it once more.
\end{tcolorbox}
\section[guastarsi]{guastarsi — to break down; to go off (food)}

\label{lesson:IT-C410-guastarsi}



\begin{quote}
Before the new one: say the Italian for to pump up, then the Italian for the breakdown.
\end{quote}

\subsection*{You'll want to know: guastarsi}
\textbf{guastarsi} — \textquotedblleft{}to break down; to go off (food)\textquotedblright{}.

\begin{quote}
La lavatrice si è guastata ieri.
\end{quote}

\begin{quote}
\emph{The washing machine broke down yesterday.}
\end{quote}

Food too: \emph{il latte si è guastato} — the milk has gone off.

\begin{cousinweb}
\emph{guastarsi} is the reflexive of \emph{guastare} 'to spoil', which is inherited from Latin \emph{vāstāre} 'to lay waste' (EtymDB).
\end{cousinweb}

\begin{grammarlens}[title={What you've built}]
One more word for telling, explaining and describing at B1.
\end{grammarlens}

\subsection*{Your turn}
\begin{itemize}
\raggedright
  \item \emph{Say it:} \emph{guastarsi}
  \item \emph{Say it:} \emph{guastarsi}, once more
  \item \emph{Say it:} a sentence of your own with \emph{guastarsi}
  \item \emph{Recall:} say the Italian for to pump up, then the Italian for the breakdown, then say \emph{guastarsi} again
\end{itemize}

\begin{tcolorbox}[breakable,colback=teal!4,colframe=teal!35!black,title={Before you move on}]
What does guastarsi mean? (\textquotedblleft{}To break down; to go off (food)\textquotedblright{}.) Say it once more.
\end{tcolorbox}
\section[il carro attrezzi]{il carro attrezzi — the tow truck, the breakdown lorry}

\label{lesson:IT-C410-carro-attrezzi}



\begin{quote}
Before the new one: say the Italian for the breakdown, then the Italian for to break down.
\end{quote}

\subsection*{You'll want to know: il carro attrezzi}
\textbf{il carro attrezzi} — \textquotedblleft{}the tow truck, the breakdown lorry\textquotedblright{}.

\begin{quote}
Abbiamo chiamato il carro attrezzi.
\end{quote}

\begin{quote}
\emph{We called the tow truck.}
\end{quote}

Also written \emph{carroattrezzi}. \emph{Gli attrezzi} are tools.

\begin{cousinweb}
\emph{carro attrezzi} is built from \emph{carro} 'cart', inherited from Latin \emph{carrus}, a word of Gaulish origin (EtymDB), and \emph{attrezzi} 'tools', whose origin is less clear.
\end{cousinweb}

\begin{grammarlens}[title={What you've built}]
One more word for telling, explaining and describing at B1.
\end{grammarlens}

\subsection*{Your turn}
\begin{itemize}
\raggedright
  \item \emph{Say it:} \emph{il carro attrezzi}
  \item \emph{Say it:} \emph{il carro attrezzi}, once more
  \item \emph{Say it:} a sentence of your own with \emph{il carro attrezzi}
  \item \emph{Recall:} say the Italian for the breakdown, then the Italian for to break down, then say \emph{il carro attrezzi} again
\end{itemize}

\begin{tcolorbox}[breakable,colback=teal!4,colframe=teal!35!black,title={Before you move on}]
What does il carro attrezzi mean? (\textquotedblleft{}The tow truck, the breakdown lorry\textquotedblright{}.) Say it once more.
\end{tcolorbox}
\section[rimorchiare]{rimorchiare — to tow}

\label{lesson:IT-C410-rimorchiare}



\begin{quote}
Before the new one: say the Italian for to break down, then the Italian for the tow truck.
\end{quote}

\subsection*{You'll want to know: rimorchiare}
\textbf{rimorchiare} — \textquotedblleft{}to tow\textquotedblright{}.

\begin{quote}
Il carro attrezzi ci rimorchia fino al paese.
\end{quote}

\begin{quote}
\emph{The tow truck tows us to the village.}
\end{quote}

\emph{Il rimorchio} is the towing, or the vehicle being towed; in slang, \emph{rimorchiare} is to chat someone up.

\begin{cousinweb}
\emph{rimorchiare} is built on \emph{rimorchio} 'towing' (EtymDB), which is inherited from Latin \emph{remulcum} 'a towing rope' (Lewis \& Short), itself from Greek.
\end{cousinweb}

\begin{grammarlens}[title={What you've built}]
One more word for telling, explaining and describing at B1.
\end{grammarlens}

\subsection*{Your turn}
\begin{itemize}
\raggedright
  \item \emph{Say it:} \emph{rimorchiare}
  \item \emph{Say it:} \emph{rimorchiare}, once more
  \item \emph{Say it:} a sentence of your own with \emph{rimorchiare}
  \item \emph{Recall:} say the Italian for to break down, then the Italian for the tow truck, then say \emph{rimorchiare} again
\end{itemize}

\begin{tcolorbox}[breakable,colback=teal!4,colframe=teal!35!black,title={Before you move on}]
What does rimorchiare mean? (\textquotedblleft{}To tow\textquotedblright{}.) Say it once more.
\end{tcolorbox}
\section[l'officina]{l'officina — the garage (for repairs), the workshop}

\label{lesson:IT-C410-officina}



\begin{quote}
Before the new one: say the Italian for the tow truck, then the Italian for to tow.
\end{quote}

\subsection*{You'll want to know: l'officina}
\textbf{l'officina} — \textquotedblleft{}the garage (for repairs), the workshop\textquotedblright{}.

\begin{quote}
La macchina resta in officina tre giorni.
\end{quote}

\begin{quote}
\emph{The car stays at the garage for three days.}
\end{quote}

A workshop for repairs; \emph{il garage} is where you keep the car.

\begin{cousinweb}
\emph{officina} is a learned borrowing from Latin \emph{officīna} 'workshop' (EtymDB), a shortened \emph{opificīna}, from \emph{opifex} 'worker', built on \emph{opus} 'work'.
\end{cousinweb}

\begin{grammarlens}[title={What you've built}]
One more word for telling, explaining and describing at B1.
\end{grammarlens}

\subsection*{Your turn}
\begin{itemize}
\raggedright
  \item \emph{Say it:} \emph{l'officina}
  \item \emph{Say it:} \emph{l'officina}, once more
  \item \emph{Say it:} a sentence of your own with \emph{l'officina}
  \item \emph{Recall:} say the Italian for the tow truck, then the Italian for to tow, then say \emph{l'officina} again
\end{itemize}

\begin{tcolorbox}[breakable,colback=teal!4,colframe=teal!35!black,title={Before you move on}]
What does l'officina mean? (\textquotedblleft{}The garage (for repairs), the workshop\textquotedblright{}.) Say it once more.
\end{tcolorbox}
